\documentclass[tikz,border=6pt]{standalone}

\usepackage{amsmath}
\usetikzlibrary{3d,arrows.meta,calc}

% =====================================================================
%  3D schematic of the Stern--Gerlach setup
%  Oven -> S1 (slit) -> SG magnet -> S2 (circular aperture) -> Detector
%
%  All coordinates are in mm.  Transverse sizes (x,z) of the magnet are
%  the real ones; the longitudinal distances (y) are schematic and can be
%  changed in the "LAYOUT" block below.
%  Beam propagates along +y. In this version the pole pieces are rotated
%  by -90 deg about y, so the field gradient and the splitting are along x.
% =====================================================================

\begin{document}

\begin{tikzpicture}[
    % ---- projection: x toward the viewer (left-down), y along the beam, z up
    x={(-0.090cm,-0.042cm)},
    y={( 0.044cm,-0.014cm)},
    z={( 0.000cm, 0.160cm)},
    line join=round, line cap=round,
    dim/.style={{Latex[length=5pt]}-{Latex[length=5pt]}, line width=0.9pt, draw=dimcolor},
    guide/.style={dotted, line width=0.6pt, draw=dimcolor},
]

% ---------------------------------------------------------------------
% Colors
% ---------------------------------------------------------------------
\definecolor{botcolor}{RGB}{150, 150, 150}   % lower pole (trench)
\definecolor{topcolor}{RGB}{195, 195, 195}   % upper pole (arch)
\definecolor{edgecolor}{RGB}{30, 30, 30}
\definecolor{dimcolor}{RGB}{60, 120, 200}
\definecolor{axiscolor}{RGB}{160, 60, 60}      % coordinate triad (x, y, z)
\definecolor{ovencolor}{RGB}{245, 160, 80}
\definecolor{coilcolor}{RGB}{190, 115, 55}     % copper windings
\colorlet{beamcolor}{red!50!blue}
\colorlet{upbeam}{blue!70!white}
\colorlet{downbeam}{red!80!white}

% ---------------------------------------------------------------------
% LAYOUT (y positions, mm) -- schematic, change freely
% ---------------------------------------------------------------------
\pgfmathsetmacro{\lOne}{110}   % oven aperture  -> S1
\pgfmathsetmacro{\lTwo}{40}   % S1             -> magnet entrance
\pgfmathsetmacro{\lSG}{100}   % magnet length
\pgfmathsetmacro{\lThr}{45}   % magnet exit    -> S2
\pgfmathsetmacro{\lFou}{105}   % S2             -> detector

\pgfmathsetmacro{\yOven}{0}
\pgfmathsetmacro{\yOvenBack}{-35}
\pgfmathsetmacro{\ySa}{\yOven+\lOne}
\pgfmathsetmacro{\yMin}{\ySa+\lTwo}
\pgfmathsetmacro{\yMout}{\yMin+\lSG}
\pgfmathsetmacro{\ySb}{\yMout+\lThr}
\pgfmathsetmacro{\yDet}{\ySb+\lFou}

\pgfmathsetmacro{\tp}{1.5}    % half thickness of the plates S1, S2

% Plates / oven / detector transverse half-sizes
\pgfmathsetmacro{\px}{12}  \pgfmathsetmacro{\pz}{10}
\pgfmathsetmacro{\ox}{12}  \pgfmathsetmacro{\oz}{11}
% ---------------------------------------------------------------------
% Detector image
%   The PNG is 840 x 1100 px; the camera frame inside its white margin
%   spans px 32..808 (horizontal) and 32..1068 (vertical), centred.
%   Measured in the frame (fractions of the frame height, from its centre,
%   + = up):  upper peak +0.140, lower peak -0.005, dark gap +0.065.
%   With the magnet rotated, the splitting is along x: the image is turned
%   by 90 deg so its "up" direction points along -x. It is shifted so the
%   gap sits on the beam axis (x = 0), and the beams end on the two peaks.
% ---------------------------------------------------------------------
\def\imgfile{SG_img_ftotal.png}
\pgfmathsetmacro{\imgH}{40}                       % frame height (now along x, mm)
\pgfmathsetmacro{\imgW}{\imgH*(776/1036)}         % frame width  (now along z)
\pgfmathsetmacro{\imgWf}{\imgW*(840/776)}         % full PNG width  (incl. margin)
\pgfmathsetmacro{\imgHf}{\imgH*(1100/1036)}       % full PNG height (incl. margin)
\pgfmathsetmacro{\imgXc}{0.065*\imgH}             % x of the frame centre
\def\imgFlip{-1}   % -1: image not mirrored on the page; 1: mirrored
\def\imgOp{0.92}   % image opacity (<1 lets the beams show faintly behind)
\pgfmathsetmacro{\dxlo}{\imgXc-\imgH/2}           % detector far edge  (-x)
\pgfmathsetmacro{\dxhi}{\imgXc+\imgH/2}           % detector near edge (+x)
\pgfmathsetmacro{\dzlo}{-\imgW/2}                 % detector bottom
\pgfmathsetmacro{\dzhi}{\imgW/2}                  % detector top

% Apertures
% (slits are long along z, perpendicular to the deflection along x)
\pgfmathsetmacro{\ovenAx}{0.6}  \pgfmathsetmacro{\ovenAz}{4.0}  % oven slit half-sizes
\pgfmathsetmacro{\saAx}{0.8}    \pgfmathsetmacro{\saAz}{5.0}    % S1 slit half-sizes
\pgfmathsetmacro{\sbR}{3.0}                                     % S2 radius

% ---------------------------------------------------------------------
% Magnet cross-section (identical to the stand-alone magnet figure)
% ---------------------------------------------------------------------
\pgfmathsetmacro{\xmin}{-10.5}
\pgfmathsetmacro{\xmax}{ 10.5}
\pgfmathsetmacro{\zmin}{-5.0}
\pgfmathsetmacro{\aa}{2.5}
\pgfmathsetmacro{\zc}{1.3*\aa}
\pgfmathsetmacro{\redge}{1.0*\aa}
\pgfmathsetmacro{\rtrench}{1.362*\aa}
\pgfmathsetmacro{\rtrenchc}{(1.3 - 1.018)*\aa}
\pgfmathsetmacro{\lw}{1.58*\aa}
\pgfmathsetmacro{\phi}{30}
\pgfmathsetmacro{\m}{tan(\phi)}
\pgfmathsetmacro{\xledge}{\rtrench + \lw*cos(\phi)}
\pgfmathsetmacro{\zledge}{\rtrenchc + \lw*sin(\phi)}
\pgfmathsetmacro{\zup}{\zc + \m*(\xmax - \redge)}

\def\poleFaceLW{1.25pt}            % profile width on the face toward S2
\def\opTop{0.85}                  % opacity of the upper pole (1 = solid)
\def\opBot{0.85}                  % opacity of the lower pole (1 = solid)

% Coils (mm): circular, centred on the magnet (mid-length, z = 0).
\pgfmathsetmacro{\cRo}{18}        % outer radius of the winding
\pgfmathsetmacro{\cRi}{6}         % radius of the iron core
\pgfmathsetmacro{\cT}{8}          % coil thickness along x
\def\cTurns{15}                    % number of winding lines drawn
\def\cOp{0.75}                    % coil opacity (1 = solid)
\pgfmathsetmacro{\cCp}{3}         % iron core sticks out of the coil by this much

% ---------------------------------------------------------------------
% Shading of the (opaque) magnet
%   brightness = \shMin + (100-\shMin) * max(0, n.L)   [% of base color]
%   L = light direction (unit vector), n = outward face normal
% ---------------------------------------------------------------------
\def\Lx{0.15}  \def\Ly{0.15}  \def\Lz{0.40}
\def\shMin{120}
% View direction (toward the viewer), derived from the x,y,z vectors
% of the tikzpicture:  V = (y1, -x1, (y2*x1 - x2*y1)/z2).  Faces with
% n.V < 0 point away from the viewer and are not drawn.
\def\Vx{0.044}  \def\Vy{0.090}  \def\Vz{0.0194}

% ---------------------------------------------------------------------
% Beam trajectories (deflection exaggerated for clarity)
% ---------------------------------------------------------------------
\pgfmathsetmacro{\dzm}{0.5}   % |z| at magnet exit
\pgfmathsetmacro{\zdet}{0.5*(0.140+0.005)*\imgH}  % |z| at detector = measured peaks
\pgfmathsetmacro{\kk}{(\zdet-\dzm)/(\yDet-\yMout)}
\pgfmathsetmacro{\zSb}{\dzm+\kk*(\ySb+\tp-\yMout)}

% Dimension line: runs along y at (x = \xdim, z = \zdim), below all parts.
% Every guide drops straight down (along -z) from the bottom edge of its
% part at x = \xdim, so all guides are parallel and start on the part.
% \xdim must lie within the x-range of every part (|x| <= 10.5 mm).
\pgfmathsetmacro{\xdim}{\xmax}
\pgfmathsetmacro{\zdim}{min(\dzlo,\xmin,-\cRo)-6}   % below the lowest part (incl. coils)

% =====================================================================
% Helper macros
% =====================================================================

% Magnet cross-section outlines at a given y  (#1 = y)
\newcommand{\topProfile}[1]{%
    (\xmin,#1,\zup) -- (\xmax,#1,\zup) -- (\redge,#1,\zc)
    -- plot[domain=360:180, samples=60, variable=\ang]
        ({\redge*cos(\ang)},#1,{\zc+\redge*sin(\ang)})
    -- cycle}
\newcommand{\botProfile}[1]{%
    (\xmin,#1,\zmin) -- (\xmax,#1,\zmin) -- (\xmax,#1,\zledge)
    -- (\xledge,#1,\zledge) -- (\rtrench,#1,\rtrenchc)
    -- plot[domain=0:180, samples=60, variable=\ang]
        ({\rtrench*cos(\ang)},#1,{\rtrenchc-\rtrench*sin(\ang)})
    -- (-\xledge,#1,\zledge) -- (\xmin,#1,\zledge) -- cycle}

% Inner (pole-face) profile lines of one end face  (#1 = y)
% Optional #1 = line width of the pole-face profile (default 1.0pt);
% the outer outline is always drawn 0.2pt thicker.
\newcommand{\poleLines}[2][1.0pt]{%
    \draw[topcolor!45!black, line width=#1]
        (\xmax,#2,\zup) -- (\redge,#2,\zc)
        -- plot[domain=360:180, samples=60, variable=\ang]
            ({\redge*cos(\ang)},#2,{\zc+\redge*sin(\ang)})
        -- (\xmin,#2,\zup);
    \draw[botcolor!45!black, line width=#1]
        (\xmax,#2,\zledge) -- (\xledge,#2,\zledge) -- (\rtrench,#2,\rtrenchc)
        -- plot[domain=0:180, samples=60, variable=\ang]
            ({\rtrench*cos(\ang)},#2,{\rtrenchc-\rtrench*sin(\ang)})
        -- (-\xledge,#2,\zledge) -- (\xmin,#2,\zledge);
    \draw[edgecolor, line width={#1+0.2pt}]
        (\xmin,#2,\zup) -- (\xmax,#2,\zup)
        (\xmin,#2,\zmin) -- (\xmax,#2,\zmin)
        (\xmin,#2,\zmin) -- (\xmin,#2,\zledge)
        (\xmax,#2,\zmin) -- (\xmax,#2,\zledge);
}

% Brightness \sh for a face with outward normal (#1,#2,#3)
\newcommand{\setshade}[3]{%
    \pgfmathtruncatemacro{\sh}{\shMin+(100-\shMin)*max(0,(#1)*\Lx+(#2)*\Ly+(#3)*\Lz)}}

% Flat magnet surface: #1 = base color, #2 = normal "nx,ny,nz",
% #3 = edge style, #4 = polygon
\newcommand{\magface}[4]{%
    \setshade#2
    \filldraw[fill=#1!\sh!black, #3] #4;}

% Thin plate centred at y = #1 ; #2 = aperture path in the xz canvas
\newcommand{\plate}[2]{%
    % side (+x) and top (+z) faces
    \fill[gray!55] (\px,#1-\tp,-\pz) -- (\px,#1+\tp,-\pz)
                -- (\px,#1+\tp,\pz) -- (\px,#1-\tp,\pz) -- cycle;
    \fill[gray!30] (-\px,#1-\tp,\pz) -- (\px,#1-\tp,\pz)
                -- (\px,#1+\tp,\pz) -- (-\px,#1+\tp,\pz) -- cycle;
    % downstream face with the aperture cut out
    \begin{scope}[canvas is xz plane at y=#1+\tp]
        \shade[left color=gray!10, right color=gray!45, even odd rule]
            (-\px,-\pz) rectangle (\px,\pz) #2;
        \draw[edgecolor, line width=0.5pt] (-\px,-\pz) rectangle (\px,\pz);
        \draw[edgecolor, line width=0.5pt] #2;
    \end{scope}
    \draw[edgecolor, line width=0.5pt]
        (\px,#1+\tp,-\pz) -- (\px,#1-\tp,-\pz) -- (\px,#1-\tp,\pz)
        -- (-\px,#1-\tp,\pz) -- (-\px,#1+\tp,\pz)
        (\px,#1-\tp,\pz) -- (\px,#1+\tp,\pz);
}

% ---------------------------------------------------------------------
% Electromagnet coils (schematic silhouettes)
%   A circular coil lies flat against the outer face of each pole piece,
%   with its axis along x (the gradient direction of the rotated magnet).
%   It is drawn as a stack of thin slices from #1 to #2 (x, mm), so the
%   cylinder and its silhouette come out right in this projection.
%   #3 = x of the face that is seen (gets the winding lines and the core),
%   or "none" when that face is hidden behind the pole.
%   The whole coil is one transparency group faded to \cOp.
% ---------------------------------------------------------------------
\pgfmathsetmacro{\cyc}{\yMin+0.5*\lSG}   % coil centre along y
% circle of radius #2 in the plane x = #1, centred on (y = \cyc, z = 0)
\newcommand{\cring}[2]{%
    plot[domain=0:360, samples=72, variable=\t] (#1,{\cyc+#2*cos(\t)},{#2*sin(\t)}) -- cycle}
\newcommand{\coil}[3]{%
    \begin{scope}[transparency group, opacity=\cOp]
    \pgfmathsetmacro{\xsecond}{#1+0.25}%
    \foreach \xs in {#1,\xsecond,...,#2}{%
        \fill[coilcolor!72!black] \cring{\xs}{\cRo};}
    \def\tmpface{#3}\def\tmpnone{none}%
    \ifx\tmpface\tmpnone\else
        \filldraw[fill=coilcolor, draw=coilcolor!45!black, line width=0.6pt]
            \cring{#3}{\cRo};
        % winding lines between the outer edge and the core
        \foreach \k in {1,...,\cTurns}{
            \pgfmathsetmacro{\wr}{\cRo-\k*(\cRo-\cRi)/(\cTurns+1)}
            \draw[coilcolor!55!black, line width=0.35pt] \cring{#3}{\wr};
        }
        % iron core through the coil
        \filldraw[fill=gray!55, draw=edgecolor, line width=0.6pt] \cring{#3}{\cRi};
    \fi
    \end{scope}}
% Iron core: solid cylinder of radius \cRi from x = #1 to x = #2 (#1 < #2),
% on the coil axis. #3 = x of the end face that gets an outline ("none" = no
% outline). It starts on the pole face, so it shows where the coil is centred.
\newcommand{\coilcore}[3]{%
    \pgfmathsetmacro{\xsecond}{#1+0.25}%
    \foreach \xs in {#1,\xsecond,...,#2}{%
        \fill[gray!45!black] \cring{\xs}{\cRi};}
    \def\tmpface{#3}\def\tmpnone{none}%
    \ifx\tmpface\tmpnone\else
        \filldraw[fill=gray!55, draw=edgecolor, line width=0.6pt] \cring{#3}{\cRi};
    \fi}

% =====================================================================
% 1) OVEN
% =====================================================================
\fill[ovencolor!80!black]
    (\ox,\yOvenBack,-\oz) -- (\ox,\yOven,-\oz) -- (\ox,\yOven,\oz) -- (\ox,\yOvenBack,\oz) -- cycle;
\fill[ovencolor!55]
    (-\ox,\yOvenBack,\oz) -- (\ox,\yOvenBack,\oz) -- (\ox,\yOven,\oz) -- (-\ox,\yOven,\oz) -- cycle;
\begin{scope}[canvas is xz plane at y=\yOven]
    \shade[left color=ovencolor!45, right color=ovencolor!85, even odd rule]
        (-\ox,-\oz) rectangle (\ox,\oz)
        (-\ovenAx,-\ovenAz) rectangle (\ovenAx,\ovenAz);
    \fill[black!80] (-\ovenAx,-\ovenAz) rectangle (\ovenAx,\ovenAz);
    \draw[edgecolor, line width=0.5pt] (-\ox,-\oz) rectangle (\ox,\oz);
\end{scope}
\draw[edgecolor, line width=0.5pt]
    (\ox,\yOven,-\oz) -- (\ox,\yOvenBack,-\oz) -- (\ox,\yOvenBack,\oz)
    -- (-\ox,\yOvenBack,\oz) -- (-\ox,\yOven,\oz)
    (\ox,\yOvenBack,\oz) -- (\ox,\yOven,\oz);

% beam: oven -> S1
\draw[beamcolor, line width=1.3pt] (0,\yOven,0) -- (0,\ySa+\tp,0);

% =====================================================================
% 2) PLATE S1  (rectangular slit)
% =====================================================================
\plate{\ySa}{(-\saAx,-\saAz) rectangle (\saAx,\saAz)}

% =====================================================================
% 3) Beam before the magnet
% =====================================================================
\draw[beamcolor, line width=1.3pt] (0,\ySa+\tp,0) -- (0,\yMin,0);

% =====================================================================
% 3b) FAR COIL, behind the far (arch) pole: its outer face at x = -\zup.
%     Drawn before the magnet, so the pole covers it; its windings and
%     core show through the transparent pole.
% =====================================================================
\pgfmathsetmacro{\xFarA}{-\zup-\cT}
\pgfmathsetmacro{\xFarB}{-\zup}
% iron core: from the pole face back through the coil, sticking out behind
\pgfmathsetmacro{\xFarCore}{\xFarA-\cCp}
\coilcore{\xFarCore}{\xFarB}{none}
\coil{\xFarA}{\xFarB}{\xFarB}   % face toward the viewer: seen through the pole

% =====================================================================
% 4) STERN--GERLACH MAGNET, rotated by -90 deg about the y axis
%    The cross-section code below is written in the magnet's own frame
%    (u, y, v) = (old x, y, old z), exactly as in the unrotated figure.
%    The scope maps it onto the setup frame:
%        x_setup = -v ,   z_setup = u
%    i.e. the arch (upper) pole ends up on the far side (-x), the trench
%    (lower) pole on the near side (+x), and the gradient points along x.
%    Inside the scope, the view direction V and the light L are expressed
%    in the magnet frame:  (Vu, Vy, Vv) = (Vz, Vy, -Vx), same for L.
%    Painter's order: far pole (behind the beam) -> beam -> near pole ->
%    end face toward S2.
% =====================================================================
\tikzset{
    botedge/.style={draw=botcolor!70!black, line width=1.0pt},
    topedge/.style={draw=topcolor!70!black, line width=1.0pt},
    outedge/.style={draw=edgecolor, line width=1.2pt},
}
\begin{scope}[x={(0cm,0.160cm)}, z={(0.090cm,0.042cm)}]
\def\Vx{0.0194}  \def\Vy{0.090}  \def\Vz{-0.044}    % V in the magnet frame
\edef\Lu{\Lz}  \edef\Lv{\Lx}
\def\Lx{\Lu}   \def\Lz{-\Lv}                         % L in the magnet frame

% --- end face toward S1 (faces away from the viewer; drawn first, it only
%     shows through the transparent poles)
\setshade{0}{-1}{0}
\begin{scope}[transparency group, opacity=\opBot]
    \fill[botcolor!\sh!black] \botProfile{\yMin};
\end{scope}
\begin{scope}[transparency group, opacity=\opTop]
    \fill[topcolor!\sh!black] \topProfile{\yMin};
\end{scope}
\pgfmathsetmacro{\opEnd}{min(\opTop,\opBot)}
\begin{scope}[transparency group, opacity=\opEnd]
    \poleLines{\yMin}
\end{scope}

% Each block below is a transparency group: it is drawn opaque and then
% faded as a whole with \opBot / \opTop, so overlapping strips and edges
% inside a pole do not produce darker seams.

\begin{scope}[transparency group, opacity=\opTop]   % ---- upper (arch) pole: far side, behind the beam
% --- both flank undersides now face the viewer
\magface{topcolor}{{0.5}{0}{-0.866}}{topedge}
    {(\redge,\yMin,\zc) -- (\xmax,\yMin,\zup)
     -- (\xmax,\yMout,\zup) -- (\redge,\yMout,\zc) -- cycle}
\magface{topcolor}{{-0.5}{0}{-0.866}}{topedge}
    {(-\redge,\yMin,\zc) -- (\xmin,\yMin,\zup)
     -- (\xmin,\yMout,\zup) -- (-\redge,\yMout,\zc) -- cycle}
% --- rounded tip, only strips facing the viewer
\foreach \i in {0,1,...,35}{
    \pgfmathsetmacro{\angA}{180+\i*5}
    \pgfmathsetmacro{\angB}{180+(\i+1)*5}
    \pgfmathsetmacro{\nx}{cos(180+(\i+0.5)*5)}
    \pgfmathsetmacro{\nz}{sin(180+(\i+0.5)*5)}
    \pgfmathsetmacro{\vis}{\nx*\Vx+\nz*\Vz}
    \ifdim\vis pt>0pt
        \setshade{\nx}{0}{\nz}
        \filldraw[topcolor!\sh!black, line width=0.3pt]
            ({\redge*cos(\angA)},\yMin,{\zc+\redge*sin(\angA)})
            -- ({\redge*cos(\angA)},\yMout,{\zc+\redge*sin(\angA)})
            -- ({\redge*cos(\angB)},\yMout,{\zc+\redge*sin(\angB)})
            -- ({\redge*cos(\angB)},\yMin,{\zc+\redge*sin(\angB)}) -- cycle;
        \draw[topedge]
            ({\redge*cos(\angA)},\yMin,{\zc+\redge*sin(\angA)})
            -- ({\redge*cos(\angB)},\yMin,{\zc+\redge*sin(\angB)});
    \fi
}
\end{scope}

% --- beam inside the magnet (deflected along v = -x_setup)
\draw[upbeam, line width=1.3pt]
    plot[domain=\yMin:\yMout, samples=40, variable=\t]
        (0,\t,{\dzm*((\t-\yMin)/\lSG)^2});
\draw[downbeam, line width=1.3pt]
    plot[domain=\yMin:\yMout, samples=40, variable=\t]
        (0,\t,{-\dzm*((\t-\yMin)/\lSG)^2});

\begin{scope}[transparency group, opacity=\opBot]   % ---- lower (trench) pole: near side, in front of the beam
% --- trench: only strips facing the viewer
\foreach \i in {0,1,...,35}{
    \pgfmathsetmacro{\angA}{\i*5}
    \pgfmathsetmacro{\angB}{(\i+1)*5}
    \pgfmathsetmacro{\nx}{-cos((\i+0.5)*5)}
    \pgfmathsetmacro{\nz}{ sin((\i+0.5)*5)}
    \pgfmathsetmacro{\vis}{\nx*\Vx+\nz*\Vz}
    \ifdim\vis pt>0pt
        \setshade{\nx}{0}{\nz}
        \filldraw[botcolor!\sh!black, line width=0.3pt]
            ({\rtrench*cos(\angA)},\yMin,{\rtrenchc-\rtrench*sin(\angA)})
            -- ({\rtrench*cos(\angA)},\yMout,{\rtrenchc-\rtrench*sin(\angA)})
            -- ({\rtrench*cos(\angB)},\yMout,{\rtrenchc-\rtrench*sin(\angB)})
            -- ({\rtrench*cos(\angB)},\yMin,{\rtrenchc-\rtrench*sin(\angB)}) -- cycle;
        \draw[botedge]
            ({\rtrench*cos(\angA)},\yMin,{\rtrenchc-\rtrench*sin(\angA)})
            -- ({\rtrench*cos(\angB)},\yMin,{\rtrenchc-\rtrench*sin(\angB)});
    \fi
}
% --- old +u side (now the top of the pole)
\magface{botcolor}{{1}{0}{0}}{outedge}
    {(\xmax,\yMin,\zmin) -- (\xmax,\yMout,\zmin)
     -- (\xmax,\yMout,\zledge) -- (\xmax,\yMin,\zledge) -- cycle}
% --- old bottom face (now the near outer face, toward the viewer)
\magface{botcolor}{{0}{0}{-1}}{outedge}
    {(\xmin,\yMin,\zmin) -- (\xmax,\yMin,\zmin)
     -- (\xmax,\yMout,\zmin) -- (\xmin,\yMout,\zmin) -- cycle}
\end{scope}

% --- end face toward S2 (drawn last, on top of everything)
\setshade{0}{1}{0}
\fill[topcolor!\sh!black, fill opacity=\opTop] \topProfile{\yMout};
\fill[botcolor!\sh!black, fill opacity=\opBot] \botProfile{\yMout};
\poleLines[\poleFaceLW]{\yMout}   % outline kept fully opaque
\end{scope}

% =====================================================================
% 4b) NEAR COIL, on the near (trench) pole: its inner face at x = -\zmin.
%     Drawn after the magnet; its outer face (toward the viewer) shows
%     the windings and the iron core.
% =====================================================================
\pgfmathsetmacro{\xNearA}{-\zmin}
\pgfmathsetmacro{\xNearB}{-\zmin+\cT}
% iron core: starts on the pole face (x = \xNearA), runs through the coil
% (seen through the transparent winding) and sticks out of its outer face
\pgfmathsetmacro{\xNearCore}{\xNearB+\cCp}
\coilcore{\xNearA}{\xNearB}{none}
\coil{\xNearA}{\xNearB}{\xNearB}
\coilcore{\xNearB}{\xNearCore}{\xNearCore}

% --- beams from the magnet exit to S2
\draw[gray, dashed, line width=0.5pt] (0,\yMout,0) -- (0,\ySb+\tp,0);
% (smooth join between the in-magnet parabola and the straight flight)
\pgfmathsetmacro{\hh}{(\ySb+\tp-\yMout)/3}
\pgfmathsetmacro{\sx}{2*\dzm/\lSG}   % slope at the magnet exit
% (the beams now split along x: "up" beam toward -x, "down" beam toward +x)
\draw[upbeam, line width=1.3pt] ({-\dzm},\yMout,0)
    .. controls ({-\dzm-\sx*\hh},\yMout+\hh,0) and ({-\zSb+\kk*\hh},\ySb+\tp-\hh,0)
    .. ({-\zSb},\ySb+\tp,0);
\draw[downbeam, line width=1.3pt] (\dzm,\yMout,0)
    .. controls ({\dzm+\sx*\hh},\yMout+\hh,0) and ({\zSb-\kk*\hh},\ySb+\tp-\hh,0)
    .. (\zSb,\ySb+\tp,0);

% gradient direction
%\draw[-{Latex[length=5pt]}, line width=0.9pt]
%    (-\xmax,\yMin+0.95*\lSG,\zup+3) -- ++(0,0,6)
%    node[above, font=\small] {$\nabla B \parallel \hat{z}$};

% =====================================================================
% 5) PLATE S2  (circular aperture)
% =====================================================================
\plate{\ySb}{(0,0) circle (\sbR)}

% beams: S2 -> detector
\draw[gray, dashed, line width=0.5pt] (0,\ySb+\tp,0) -- (0,\yDet,0);
\draw[upbeam,   line width=1.3pt] ({-\zSb},\ySb+\tp,0) -- ({-\zdet},\yDet,0);
\draw[downbeam, line width=1.3pt] (\zSb,\ySb+\tp,0)    -- (\zdet,\yDet,0);

% =====================================================================
% 6) DETECTOR with the measured image (\imgfile)
% =====================================================================
% The image is mirrored horizontally (xscale) and then turned by 90 deg
% (rotate), so that its "up" direction points along -x, where the upper
% peak meets the "up" beam, without appearing mirrored on the page.
\begin{scope}[canvas is xz plane at y=\yDet]
    \begin{scope}
        \clip (\dxlo,\dzlo) rectangle (\dxhi,\dzhi);
        \node[transform shape, inner sep=0pt, anchor=center, opacity=\imgOp,
              rotate=90, xscale=\imgFlip] at (\imgXc,0)
            {\includegraphics[width=\imgWf cm, height=\imgHf cm]{\imgfile}};
    \end{scope}
    \draw[dimcolor!70!black, line width=0.6pt] (\dxlo,\dzlo) rectangle (\dxhi,\dzhi);
\end{scope}

% =====================================================================
% 7) DIMENSIONS
% =====================================================================
% Guides start exactly at the near bottom corner of each part and drop
% vertically (on the page) until they meet the dimension line, so they
% stay parallel whatever the part sizes and distances are.
\coordinate (dimA) at (\xdim,\yOven,\zdim);   % two points defining
\coordinate (dimB) at (\xdim,\yDet,\zdim);    % the dimension line
% #1 = name of the foot point on the dimension line, #2 = start corner
\newcommand{\dimguide}[2]{%
    \coordinate (c#1) at #2;
    \coordinate (d#1) at ($(c#1)+(0,-1cm)$);
    \coordinate (#1) at (intersection of c#1--d#1 and dimA--dimB);
    \draw[guide] (c#1) -- (#1);}
\dimguide{g0}{(\ox,\yOven,-\oz)}      % oven aperture face
\dimguide{g1}{(\px,\ySa,-\pz)}        % S1 (centre of the plate)
\dimguide{g2}{({-\zmin},\yMin,\xmin)}    % magnet entrance (rotated magnet)
\dimguide{g3}{({-\zmin},\yMout,\xmin)}   % magnet exit
\dimguide{g4}{(\px,\ySb,-\pz)}        % S2 (centre of the plate)
\dimguide{g5}{(\dxhi,\yDet,\dzlo)}    % detector

\foreach \a/\b/\lab in {g0/g1/{\ell_1}, g1/g2/{\ell_2}, g2/g3/{\ell_{\mathrm{SG}}},
                        g3/g4/{\ell_3}, g4/g5/{\ell_4}}{
    \draw[dim] (\a) -- (\b);
    \node[below left=-1pt, font=\small] at ($(\a)!0.5!(\b)$) {$\lab$};
}

% =====================================================================
% 8) LABELS
% =====================================================================
\node[font=\bfseries\small, anchor=south] at (-\ox,{0.5*(\yOven+\yOvenBack)},\oz+1) {Oven};
\node[font=\small, anchor=south west]     at (-\px,\ySa,\pz) {$\mathcal{S}_1$};
\node[font=\small, anchor=south west]     at (-\px,\ySb,\pz) {$\mathcal{S}_2$};

% Tilted labels: the text is laid in a plane of the setup (transform shape).
% The canvas stretches 1 cm into one 3D unit (1 mm), so the text is scaled
% back by 1/|axis vector| to keep its normal size on the page:
%   |x| = 0.0993 cm, |y| = 0.0462 cm, |z| = 0.160 cm.
\pgfmathsetmacro{\sX}{1/0.0993}
\pgfmathsetmacro{\sY}{1/0.0462}
\pgfmathsetmacro{\sZ}{1/0.160}
% magnet: runs along the magnet (y), centred on it (mid-length, x = 0)
\begin{scope}[canvas is yz plane at x=0]
    \node[transform shape, xscale=\sY, yscale=\sZ, font=\bfseries\small, anchor=south]
        at ({\yMin+0.75*\lSG},\xmax+16) {Stern--Gerlach magnet};   % raised above the far coil   % top of the rotated magnet is z = \xmax
\end{scope}
% detector: runs along the detector sheet (x), just above its top edge;
% xscale is negative because +x points to the left on the page
\begin{scope}[canvas is xz plane at y=\yDet]
    \node[transform shape, xscale=-\sX, yscale=\sZ, font=\bfseries\small, anchor=south]
        at (\imgXc,\dzhi+1) {Detector};
\end{scope}

% =====================================================================
% 9) COORDINATE TRIAD
% =====================================================================
\begin{scope}[shift={(0,\yDet-25,28)}]
    % Labels follow the rotated magnet: the gradient direction is called z.
    % In drawing coordinates:  z_label = -x_drawing,  x_label = +z_drawing,
    % y unchanged.
    \draw[-{Latex[length=5pt]}, line width=1.2pt, axiscolor] (0,0,0) -- (-12,0,0) node[above right=-2pt] {$z$};
    \draw[-{Latex[length=5pt]}, line width=1.2pt, axiscolor] (0,0,0) -- (0,30,0)  node[right]             {$y$};
    \draw[-{Latex[length=5pt]}, line width=1.2pt, axiscolor] (0,0,0) -- (0,0,9)   node[above]             {$x$};
\end{scope}

\end{tikzpicture}

\end{document}
